% Section 7. Input from paper.tex; every number traces to sources.md and is asserted by check_numbers.py.
\providecommand{\lm}{\texttt{local\_mean\_28}}
\providecommand{\qs}{QScore9}
\providecommand{\hc}{\texttt{hurdle\_count\_28}}

\section{Discussion}
\label{sec:discussion}

\paragraph{What the three studies say together.} The models' value on sparse daily counts is in the size of
the positive days. Their gain over a trailing mean sits on positive days, where the local mean overpredicts on
GitHub and npm and underpredicts on NYC 311. Against count laws, their positive law beats a shifted Poisson
wherever that law is misspecified, and an empirical law on dense series. T2 rejects for the dynamic hurdle
against the empirical count hurdle in 14 of 18 registered tests on NYC 311, three of them by rank sum with a
median of exactly zero. On sparse series that contrast owes at least as much to the swapped occurrence as to
the positive law, and on very sparse series everything.

Where series are dense, the models' own probability of a zero day is at least as good as a context-only
estimate. That holds under the quantile score in all three settings, and under the occurrence Brier score on
NYC 311 and at a week on npm. This is confirmatory on NYC 311, through the dense stratum's own six tests of
the dynamic occurrence. On npm and the grid the confirmatory families are pooled, none favors the fixed
substitution and two of six npm tests and all three grid tests reject against it. Reading them as the dense
answer therefore rests on the declared secondary strata.

On sparse series the answer depends on the domain. The typical sparse npm series favors a context estimate at
a week, for the TSB hurdle as much as for the Jeffreys estimate, and not at a month by the within-stratum
test. On the sparse NYC 311 stratum one of six T1 tests favors it and four of six medians are exactly zero. At
a month, the grid's cells with $z = 0.95$ favor it as well. On the natural panels nothing about occurrence can
be read on series that are zero nine days in ten, because the completion sets the models' occurrence
probability and the quantile score ties.

\paragraph{The two readings of the accounting trade.} Study 2 left two readings of the net sign open. One is
that the context estimate is the worse occurrence forecast. The other is that the hurdle's truncated law is
not the model's size forecast. Of the two, the record supports the first on dense NYC 311 series and at a week
on the grid. It supports the second at a month on the grid and, on the Brier score, on npm, post hoc there
(Appendix~\ref{app:s3tables}). Nothing beyond that is attributed here.

\paragraph{Expected and observed.} Each study was registered with an expectation, and none was borne out in
full. Study 1's registered hypothesis predicted lower scores for the models than for both references. Against
weekly AutoETS the models did score lower, but against the local mean they won only on the mean, and the
typical repository favored the reference in every cell. Study 2 asked whether replacing the models' zero mass
with a context-only occurrence estimate would improve the score by removing the small zero-day loss. Pooled,
it did not. Study 3's registered hypothesis predicted that the sparse npm reading would reproduce on a second
domain. On the sparse and very sparse strata it was not confirmed, and three features of the record explain
why.

The first is the completion floor, which was in the registered rule from the start. Before the Study 3 run,
its first stage counted 74 of 100 very sparse series on the Jeffreys fallback, and its record anticipated many
exact ties. On a series that is zero nine days in ten the registered ``model's own zero mass'' is a constant
between 0.0607 and 0.1. Thus, the very sparse question is not one the registered construction can ask, and the
ties were foreseeable from the definition. The second is that the sparse npm week appears to be a property of
the context estimate as such, since the TSB hurdle wins the same series. It is also a property of the week
rather than the month.

The third is the registered dynamic count, which beats every foundation model on the grid at a month and loses
to them on npm and NYC 311, on the mean and by count. With the occurrence held equal, the foundation models'
positive laws underpredict the grid's positive days more than its negative binomial does
(Appendix~\ref{app:s2tables}). One further design choice limits the reading. The threshold of 0.5 between the
regimes was inherited from Study 2's sampling strata, not derived. Figure~\ref{fig:signature} suggests that
the point at which the substitution starts to help is domain-specific. If so, a stratum boundary fixed in
advance may obscure where that point lies on a given domain.

\paragraph{Implications for a follow-up study.} A follow-up study of this question would register the
occurrence log or Brier score and the discrete ranked probability score as its primaries. Its model output
would be one that can express an occurrence probability below one in ten, such as Chronos-2's own 21 quantile
levels, a mixture or sample paths. TSB and iETS would be registered beside the count baselines, the regime
threshold would be chosen from a per-series curve, and the study would run on a prospective window. Together
these choices define a single study, plausibly the one a reader of the intermittent-demand literature would
ask for. Short of such a study, a reviewer could ask for a registered replication of the sparse npm week on
target days not yet read, six tests under the occurrence Brier score. That replication is held in reserve.

\paragraph{Scope of the claim.} The claim covers the registered combination and nothing broader. That
combination is a context-only or dynamic occurrence probability substituted for the zero mass of frozen
Chronos-2 and TimesFM forecasts, their completed positive law kept, judged by a proper quantile score against
a completed-unwrapped control and matched count hurdles, with a dynamic count baseline beside them. It is
tested on a hash-sampled panel of daily npm download counts, on the registered synthetic grid and, stratum by
stratum, on a stratified panel of daily NYC 311 service-request counts, with the mean and the typical series
reported side by side. Study 1's GitHub star counts enter through the raw comparison alone. Within that scope
the answer is the one the title gives, size and not occurrence.
